\begin{abstract}
  Scene text in video---signs, labels, menus, and printed matter
  captured on camera---is a natural target for automated cross-language
  localization, but end-to-end generative video editors currently lack
  the glyph-level control needed to reliably render a specific target
  string. We describe a prototype system that takes the opposite
  approach: a modular five-stage pipeline that explicitly detects,
  tracks, edits, and composites scene text, preserving control over the
  target string at every step. The system integrates seven
  machine-learning models, of which two we implemented and trained
  ourselves: an Alignment Refiner of our own design that addresses
  residual tracking error, and a compact Blur Prediction Network
  that follows STRIVE's per-frame blur formulation but is
  re-implemented from scratch (STRIVE's code is unreleased) with
  our own backbone, training data, and inference-time fixes. We describe the design of each stage, the
  engineering decisions needed to handle long-to-short translations
  and interchangeable background inpainting, and we evaluate on real
  video examples. Results show the feasibility of the approach and
  position the system as a reproducible open baseline in the
  glyph-controlled paradigm, complementary to generative video editors
  that currently offer superior style preservation at the cost of
  glyph accuracy.
\end{abstract}
